% Chapter 3 (Weeks 5--8) — Linear Transformations
% Sources for Weeks 7--8: :contentReference[oaicite:0]{index=0} :contentReference[oaicite:1]{index=1} :contentReference[oaicite:2]{index=2} :contentReference[oaicite:3]{index=3}
%====================================================================

\section{Linear Transformations}

%====================================================================

\subsection{Definitions, linearity tests, and core examples (Week 5)}

\begin{definition}{Linear transformation}{def:linear-transformation}
Let $V$ and $W$ be vector spaces over the same field $\mathbb{F}$.
A function $T:V\to W$ is a \emph{linear transformation} if for all $u,v\in V$ and all $\lambda\in\mathbb{F}$,
\[
T(u+v)=T(u)+T(v)
\quad\text{and}\quad
T(\lambda u)=\lambda\,T(u).
\]
Equivalently, $T$ is linear if and only if for all $u,v\in V$ and all $\alpha,\beta\in\mathbb{F}$,
\[
T(\alpha u+\beta v)=\alpha\,T(u)+\beta\,T(v).
\]
\end{definition}

Intuitively, a linear transformation is a ``structure-preserving'' map: it respects the two operations that define a vector space (addition and scalar multiplication).
An equivalent one-shot test is to check
\[
\forall v_1,\dots,v_k\in V,\ \forall \lambda_1,\dots,\lambda_k\in\mathbb{F}:\qquad
T\!\left(\sum_{i=1}^k \lambda_i v_i\right)=\sum_{i=1}^k \lambda_i\,T(v_i),
\]
i.e.\ $T$ commutes with \emph{all} finite linear combinations.
This general form follows by induction from the two-vector case and is often what you actually use in proofs.
Two immediate consequences that serve as quick sanity checks:
\begin{itemize}
\item $T(0_V)=0_W$ always (set $\lambda=0$ in homogeneity). If $T(0)\neq 0$, then $T$ is \emph{not} linear.
\item $T(-v)=-T(v)$ for all $v\in V$ (set $\lambda=-1$).
\end{itemize}

\begin{keyformula}
\textbf{Linearity test (one line).}
To prove $T:V\to W$ is linear, show
\[
\forall u,v\in V,\ \forall \alpha,\beta\in\mathbb{F}:\qquad
T(\alpha u+\beta v)=\alpha\,T(u)+\beta\,T(v).
\]
To disprove linearity, it suffices to find one counterexample to additivity or homogeneity.
\end{keyformula}

\begin{definition}{The vector space of linear maps}{def:space-of-linear-maps}
Let $V$ and $W$ be vector spaces over $\mathbb{F}$. Define
\[
\mathcal{L}(V,W)\coloneqq \{T:V\to W \mid T\text{ is linear}\},
\qquad
\mathcal{L}(V)\coloneqq \mathcal{L}(V,V).
\]
For $S,T\in\mathcal{L}(V,W)$ and $\lambda\in\mathbb{F}$, define pointwise operations:
\[
(S+T)(v)\coloneqq S(v)+T(v),\qquad (\lambda T)(v)\coloneqq \lambda\,T(v)\qquad(\forall v\in V).
\]
\end{definition}

\begin{theorem}{$\mathcal{L}(V,W)$ is a vector space}{thm:LVW-vector-space}
Let $V,W$ be vector spaces over $\mathbb{F}$. Then $\mathcal{L}(V,W)$ is a vector space over $\mathbb{F}$
under pointwise addition and scalar multiplication.
\end{theorem}

\begin{proof}
We verify closure; the remaining axioms follow pointwise from those in $W$.

Let $S,T\in\mathcal{L}(V,W)$ and $\lambda\in\mathbb{F}$. For all $u,v\in V$,
\[
(S+T)(u+v)=S(u+v)+T(u+v)=S(u)+S(v)+T(u)+T(v)=(S+T)(u)+(S+T)(v),
\]
and similarly $(S+T)(\lambda u)=\lambda(S+T)(u)$, hence $S+T\in\mathcal{L}(V,W)$.
Also,
\[
(\lambda T)(u+v)=\lambda T(u+v)=\lambda T(u)+\lambda T(v)=(\lambda T)(u)+(\lambda T)(v),
\qquad
(\lambda T)(\mu u)=\mu(\lambda T)(u),
\]
so $\lambda T\in\mathcal{L}(V,W)$.
\end{proof}

When both $V$ and $W$ are finite-dimensional with $\dim(V)=m$ and $\dim(W)=n$, we will later show that $\mathcal{L}(V,W)\cong \mathbb{F}^{n\times m}$ as vector spaces (see the extension box in Section~\ref{def:change-of-coordinates}ff.), hence
\[
\dim(\mathcal{L}(V,W))=mn.
\]
This matches the dimension of $\mathbb{F}^{n\times m}$ and reinforces the idea that choosing bases turns abstract linear maps into concrete matrices.

\begin{examtip}
\textbf{Tutorial 5, Question 1 note.}
When asked to show $\mathcal{L}(V,W)$ is a vector space, you can either
(i) verify all eight axioms directly (pointwise, borrowing from $W$), or
(ii) show $\mathcal{L}(V,W)$ is a subspace of the function space $W^V$ by checking closure under addition and scalar multiplication plus containing the zero map. Method (ii) is usually faster on exams because the axioms are inherited automatically.
\end{examtip}

\begin{examplebox}
\textbf{Worked Example: linear vs.\ non-linear maps on $\mathbb{R}^2$.}
Let $T_1,T_2:\mathbb{R}^2\to\mathbb{R}^2$ be defined by
\[
T_1(x,y)=(x+2y,\,3x-y),
\qquad
T_2(x,y)=(x^2,\,y).
\]
\begin{itemize}
\item $T_1$ is linear: both additivity and homogeneity follow by direct expansion.
\item $T_2$ is not linear: homogeneity fails since $T_2(2,0)=(4,0)\neq 2(1,0)=2T_2(1,0)$.
\end{itemize}
\end{examplebox}

\begin{examplebox}
\textbf{Worked Example: common linear maps (from Week~5 lectures).}
\begin{itemize}
\item \textbf{Derivative operator:} $D:\mathcal{P}(\mathbb{F})\to\mathcal{P}(\mathbb{F})$, $D(P)=P'$, is linear because $(P+Q)'=P'+Q'$ and $(\lambda P)'=\lambda P'$.
\item \textbf{Multiplication operator:} Fix $Q\in\mathcal{P}(\mathbb{F})$; then $M_Q:\mathcal{P}(\mathbb{F})\to\mathcal{P}(\mathbb{F})$, $M_Q(P)=QP$, is linear because polynomial multiplication distributes over addition and commutes with scalar multiplication.
\item \textbf{Matrix operator:} Fix $A\in\mathbb{F}^{m\times n}$; then $T_A:\mathbb{F}^{n\times 1}\to\mathbb{F}^{m\times 1}$, $T_A(x)=Ax$, is linear by properties of matrix multiplication.
\item \textbf{Backward-shift:} On $\mathbb{F}^\infty$, $\sigma_\leftarrow(x_1,x_2,x_3,\dots)=(x_2,x_3,x_4,\dots)$ is linear (check componentwise).
\end{itemize}
These examples illustrate how linearity appears across many settings: polynomials, matrices, and sequence spaces.
\end{examplebox}

%====================================================================

\subsection{Kernel (null space) and range (image) (Week 5)}

\begin{definition}{Null space and range}{def:null-range}
Let $V,W$ be vector spaces over $\mathbb{F}$ and let $T\in\mathcal{L}(V,W)$.
\[
\begin{aligned}
\operatorname{Null}(T)&\coloneqq \{v\in V \mid T(v)=0\}
\quad(\text{also written }\ker(T)),\\
\operatorname{Range}(T)&\coloneqq \{T(v)\in W \mid v\in V\}
\quad(\text{also written }\operatorname{im}(T)).
\end{aligned}
\]
\end{definition}

The null space measures ``how far $T$ is from being injective'': elements of $\operatorname{Null}(T)$ are the vectors that $T$ cannot distinguish from $0$.
The range measures ``how much of $W$ does $T$ reach'': it is the set of all possible outputs.
Together with the Rank--Nullity Theorem (next subsection), these two subspaces completely determine the ``information loss'' and ``coverage'' of $T$.

\begin{theorem}{Null space and range are subspaces}{thm:null-range-subspaces}
Let $V,W$ be vector spaces over $\mathbb{F}$ and $T\in\mathcal{L}(V,W)$.
Then $\operatorname{Null}(T)$ is a subspace of $V$, and $\operatorname{Range}(T)$ is a subspace of $W$.
\end{theorem}

\begin{proof}
\textbf{Null space.} Since $T(0)=0$, we have $0\in\operatorname{Null}(T)$.
If $v,\tilde v\in\operatorname{Null}(T)$ and $\lambda\in\mathbb{F}$, then
\[
T(v+\tilde v)=T(v)+T(\tilde v)=0+0=0,\qquad T(\lambda v)=\lambda T(v)=\lambda\cdot 0=0,
\]
so $v+\tilde v,\lambda v\in\operatorname{Null}(T)$.

\textbf{Range.} Since $0=T(0)\in\operatorname{Range}(T)$, it is nonempty.
If $w=T(v)$ and $w'=T(v')$ are in the range, then for $\lambda\in\mathbb{F}$,
\[
w+w'=T(v)+T(v')=T(v+v')\in\operatorname{Range}(T),\qquad
\lambda w=\lambda T(v)=T(\lambda v)\in\operatorname{Range}(T).
\]
Thus $\operatorname{Range}(T)$ is a subspace of $W$.
\end{proof}

\begin{extension}[title={Extension (Beyond Syllabus: Preimages preserve subspaces)}]
\textbf{Syllabus link:} MH1201 Weeks 5--6 — Kernel and range.

\textbf{Lemma.} Let $T:V\to W$ be linear and let $U\le W$ be a subspace. Then
\[
T^{-1}(U)\coloneqq \{v\in V \mid T(v)\in U\}
\]
is a subspace of $V$.

\textbf{Proof.} Since $0\in U$ and $T(0)=0$, we have $0\in T^{-1}(U)$.
Let $v,\tilde v\in T^{-1}(U)$ and $\lambda\in\mathbb{F}$. Then $T(v),T(\tilde v)\in U$.
Because $U$ is a subspace, $T(v)+T(\tilde v)\in U$ and $\lambda T(v)\in U$.
By linearity,
\[
T(v+\tilde v)=T(v)+T(\tilde v)\in U,\qquad T(\lambda v)=\lambda T(v)\in U,
\]
so $v+\tilde v,\lambda v\in T^{-1}(U)$. Hence $T^{-1}(U)$ is a subspace. \qed
\end{extension}

\begin{examplebox}
\textbf{Worked Example: kernel and range of a polynomial map.}
Let $T:\mathcal{P}_2(\mathbb{R})\to\mathbb{R}^2$ be defined by
\[
\forall a,b,c\in\mathbb{R}:\qquad T(a+bX+cX^2)=(a+c,\ a-b).
\]
\textbf{Kernel.} Solve $a+c=0$ and $a-b=0$, giving $(a,b,c)=(a,a,-a)$.
Thus
\[
\operatorname{Null}(T)=\{a(1+X-X^2)\mid a\in\mathbb{R}\}=\operatorname{span}\{1+X-X^2\}.
\]
\textbf{Range.} For any $(u,v)\in\mathbb{R}^2$, take $a=0$, $b=-v$, $c=u$, giving $T(uX^2-vX)=(u,v)$.
Hence $\operatorname{Range}(T)=\mathbb{R}^2$.
\end{examplebox}

\begin{examplebox}
\textbf{Worked Example (Tutorial 5, Question 3 style): kernel, range, and non-direct-sum.}
Define $T:\mathbb{R}^{2,2}\to\mathbb{R}^{2,2}$ by $T(A)=AJ-\operatorname{tr}(A)\,J$ where $J=\begin{bmatrix}1&1\\1&1\end{bmatrix}$.
Explicit computation on $A=\begin{bmatrix}a&b\\c&d\end{bmatrix}$ gives
\[
T(A)=\begin{bmatrix}b-d&b-d\\c-a&c-a\end{bmatrix}.
\]
\textbf{Range:} Every output has equal columns, and any such matrix is attainable. Thus
$\operatorname{Range}(T)=\Bigl\{\begin{bmatrix}u&u\\v&v\end{bmatrix}\Bigr\}$, with $\operatorname{Rank}(T)=2$.
\textbf{Kernel:} $T(A)=0$ iff $b=d$ and $c=a$, so
$\operatorname{Null}(T)=\Bigl\{\begin{bmatrix}a&b\\a&b\end{bmatrix}\Bigr\}$, with $\operatorname{Nullity}(T)=2$.
\textbf{Sanity check:} $\operatorname{Rank}+\operatorname{Nullity}=2+2=4=\dim(\mathbb{R}^{2,2})$. \checkmark

\textbf{Non-direct-sum:} $J=\begin{bmatrix}1&1\\1&1\end{bmatrix}$ lies in both $\operatorname{Null}(T)$ and $\operatorname{Range}(T)$, so $\operatorname{Null}(T)\cap\operatorname{Range}(T)\neq\{0\}$ and the sum is \emph{not} direct.
This shows that $\mathbb{R}^{2,2}=\operatorname{Null}(T)\oplus\operatorname{Range}(T)$ does \emph{not} hold in general.
\end{examplebox}

%====================================================================

\subsection{Rank--Nullity theorem and injective/surjective criteria (Weeks 5--6)}

\begin{definition}{Rank and nullity}{def:rank-nullity-def}
Let $V,W$ be vector spaces over $\mathbb{F}$ with $V$ finite-dimensional, and let $T\in\mathcal{L}(V,W)$.
Define
\[
\operatorname{Nullity}(T)\coloneqq \dim(\operatorname{Null}(T)),
\qquad
\operatorname{Rank}(T)\coloneqq \dim(\operatorname{Range}(T)).
\]
\end{definition}

\begin{theorem}{Rank--Nullity Theorem}{thm:rank-nullity}
Let $V,W$ be vector spaces over $\mathbb{F}$ with $V$ finite-dimensional, and let $T\in\mathcal{L}(V,W)$.
Then
\[
\dim(V)=\operatorname{Rank}(T)+\operatorname{Nullity}(T).
\]
\end{theorem}

\begin{proof}
Let $m\coloneqq \dim(\operatorname{Null}(T))$ and choose a basis $(b_1,\dots,b_m)$ of $\operatorname{Null}(T)$.
Extend it to a basis of $V$:
\[
(b_1,\dots,b_m,v_1,\dots,v_{n-m}),
\qquad n\coloneqq \dim(V).
\]
We claim $(T(v_1),\dots,T(v_{n-m}))$ is a basis of $\operatorname{Range}(T)$.

\textbf{Independence.} If $\sum_{i=1}^{n-m}\lambda_i T(v_i)=0$, then
$T(\sum_{i=1}^{n-m}\lambda_i v_i)=0$, so $\sum_{i=1}^{n-m}\lambda_i v_i\in\operatorname{Null}(T)$.
But it also lies in $\operatorname{span}\{v_1,\dots,v_{n-m}\}$; by uniqueness of coordinates in the displayed basis, this forces
$\sum_{i=1}^{n-m}\lambda_i v_i=0$, hence all $\lambda_i=0$.

\textbf{Spanning.} If $w\in\operatorname{Range}(T)$, then $w=T(v)$ for some $v\in V$.
Write $v=\sum_{j=1}^m\alpha_j b_j+\sum_{i=1}^{n-m}\beta_i v_i$. Since $T(b_j)=0$,
\[
w=T(v)=\sum_{i=1}^{n-m}\beta_i T(v_i)\in \operatorname{span}\{T(v_1),\dots,T(v_{n-m})\}.
\]
Thus $\operatorname{Rank}(T)=n-m$ and $\operatorname{Nullity}(T)=m$, giving $n=(n-m)+m$.

\end{proof}

The Rank--Nullity Theorem is one of the most frequently used results in linear algebra.
It says that the ``dimensions consumed by the kernel'' plus the ``dimensions captured in the range'' always add up to the dimension of the domain.

\begin{examtip}
\textbf{Using Rank--Nullity on exams.}
\begin{itemize}
\item If you know any two of $\{\dim(V),\operatorname{Rank}(T),\operatorname{Nullity}(T)\}$, you can deduce the third.
\item To show a map is injective, it suffices to show $\operatorname{Nullity}(T)=0$. To show surjectivity onto $W$, show $\operatorname{Rank}(T)=\dim(W)$.
\item Always state the hypothesis ``$V$ is finite-dimensional'' before applying the theorem.
\item Sanity check: after computing kernel and range independently, verify $\operatorname{Rank}+\operatorname{Nullity}=\dim(V)$.
\end{itemize}
\end{examtip}

\begin{definition}{Injective, surjective, bijective}{def:inj-surj-bij}
Let $T\in\mathcal{L}(V,W)$.
\begin{itemize}
\item $T$ is \emph{injective} if $T(v)=T(\tilde v)\Rightarrow v=\tilde v$.
\item $T$ is \emph{surjective onto $W$} if $\forall w\in W,\ \exists v\in V:\ T(v)=w$.
\item $T$ is \emph{bijective} if it is both injective and surjective.
\end{itemize}

\end{definition}

Injectivity means distinct inputs yield distinct outputs (no information loss); surjectivity means every element of the codomain is reached (full coverage).
For linear maps, the null-space test (Theorem~\ref{thm:null-space-test-injective}) provides a more efficient check for injectivity: instead of comparing all pairs, you only need to check that the kernel is trivial.

\begin{keyformula}
\textbf{RREF criteria for injectivity/surjectivity (Tutorial 6, Question 2).}
Let $V,W$ be finite-dimensional with ordered bases $\beta,\gamma$. Let $T\in\mathcal{L}(V,W)$ and $A=[T]^{\gamma}_{\beta}\in\mathbb{F}^{n\times m}$.
\begin{itemize}
\item $T$ injective $\Longleftrightarrow$ RREF of $A$ has a pivot in every column ($\Leftrightarrow \operatorname{Rank}(A)=m$, no free variables).
\item $T$ surjective onto $W$ $\Longleftrightarrow$ RREF of $A$ has a pivot in every row ($\Leftrightarrow \operatorname{Rank}(A)=n$).
\end{itemize}
These translate injectivity/surjectivity into concrete row-reduction checks that are fast to perform on exams.
\end{keyformula}

\begin{keyformula}
\textbf{Equal-dimension criterion (Tutorial 6, Question 3).}
If $\dim(V)=\dim(W)=n$ and $T\in\mathcal{L}(V,W)$, then
\[
T\text{ injective}\ \Longleftrightarrow\ T\text{ surjective onto }W\ \Longleftrightarrow\ T\text{ bijective}\ \Longleftrightarrow\ [T]^{\gamma}_{\beta}\text{ invertible}.
\]
This is because Rank--Nullity forces $\operatorname{Rank}(T)=n-\operatorname{Nullity}(T)$, so nullity $0$ is equivalent to rank $n=\dim(W)$.
\end{keyformula}

\begin{theorem}{Null-space test for injectivity}{null-space-test-injective}
Let $T\in\mathcal{L}(V,W)$. Then $T$ is injective if and only if $\operatorname{Null}(T)=\{0\}$.
\end{theorem}

\begin{proof}
If $T$ is injective, then $T(v)=0=T(0)$ implies $v=0$, so $\operatorname{Null}(T)=\{0\}$.
Conversely, if $\operatorname{Null}(T)=\{0\}$ and $T(v)=T(\tilde v)$, then $T(v-\tilde v)=0$, hence $v-\tilde v=0$ and $v=\tilde v$.
\end{proof}

\begin{keyformula}
\textbf{Finite-dimensional shortcuts (state hypotheses!).}
Assume $V,W$ are finite-dimensional.
\begin{itemize}
\item $T$ injective $\Longleftrightarrow \operatorname{Nullity}(T)=0 \Longleftrightarrow \operatorname{Rank}(T)=\dim(V)$.
\item $T$ surjective onto $W$ $\Longleftrightarrow \operatorname{Rank}(T)=\dim(W)$.
\item If $\dim(V)=\dim(W)$, then injective $\Longleftrightarrow$ surjective $\Longleftrightarrow$ bijective.
\end{itemize}
\end{keyformula}

\begin{definition}{Invertible linear transformation}{def:invertible-linear-map}
Let $T\in\mathcal{L}(V,W)$. We say $T$ is \emph{invertible} if there exists a function $S:W\to V$ such that
\[
S\circ T=I_V\quad\text{and}\quad T\circ S=I_W.
\]
In that case, $S$ is unique and is denoted by $T^{-1}$.
\end{definition}

\begin{theorem}{Invertibility $\Longleftrightarrow$ bijectivity; inverse is linear}{thm:invertible-bijective}
Let $T\in\mathcal{L}(V,W)$. Then $T$ is invertible if and only if $T$ is bijective.
Moreover, if $T$ is invertible, then $T^{-1}\in\mathcal{L}(W,V)$.
\end{theorem}

\begin{proof}
If $T$ has inverse $T^{-1}$, then injectivity and surjectivity are immediate from $T^{-1}\circ T=I_V$ and $T\circ T^{-1}=I_W$.
Conversely, if $T$ is bijective, define $T^{-1}(w)$ to be the unique $v$ with $T(v)=w$; then $T^{-1}\circ T=I_V$ and $T\circ T^{-1}=I_W$.

To show $T^{-1}$ is linear: let $w_1,w_2\in W$ and set $v_i=T^{-1}(w_i)$. Then $T(v_1+v_2)=w_1+w_2$, so
$T^{-1}(w_1+w_2)=v_1+v_2$, and similarly $T^{-1}(\lambda w_1)=\lambda v_1$. Hence $T^{-1}$ is linear.

\end{proof}

\begin{keyformula}
\textbf{Composition, injectivity, and surjectivity (Tutorial 5, Questions 6--7; Tutorial 6, Question 4).}
Let $U,V,W$ be vector spaces and $S\in\mathcal{L}(U,V)$, $T\in\mathcal{L}(V,W)$.
\begin{itemize}
\item If $T\circ S$ is injective, then $S$ is injective (but $T$ need not be).
\item If $T\circ S$ is surjective onto $W$, then $T$ is surjective onto $W$ (but $S$ need not be).
\item If $S$ and $T$ are both invertible, then $T\circ S$ is invertible and $(T\circ S)^{-1}=S^{-1}\circ T^{-1}$.
\item The converse fails: $T\circ S$ invertible does \emph{not} imply both $S,T$ invertible individually.
\end{itemize}
\end{keyformula}

\begin{examplebox}
\textbf{Worked Example (Tutorial 5, Question 7): counterexamples for converses.}
Take $U=\mathbb{R}$, $V=\mathbb{R}^2$, $W=\mathbb{R}$.
Define $S(t)=(t,0)$ and $T(x,y)=x$.
Then $(T\circ S)(t)=t$ is bijective.
\begin{itemize}
\item $T\circ S$ is injective, but $T$ is \emph{not} injective (since $T(0,1)=0=T(0,0)$).
\item $T\circ S$ is surjective, but $S$ is \emph{not} surjective onto $\mathbb{R}^2$ (since $(0,1)\notin S(\mathbb{R})$).
\end{itemize}
The key insight: $T\circ S$ injective only forces $\operatorname{Range}(S)\cap\ker(T)=\{0\}$, not $\ker(T)=\{0\}$.
\end{examplebox}

%====================================================================

\subsection{Matrix representations and composition (Week 6)}

\begin{definition}{Matrix of a linear map}{matrix-rep}
Let $V,W$ be finite-dimensional over $\mathbb{F}$ with ordered bases
\[
\beta=(b_1,\dots,b_n)\ \text{for }V,\qquad \gamma=(c_1,\dots,c_m)\ \text{for }W,
\]
and let $T\in\mathcal{L}(V,W)$.
Define $[T]^{\gamma}_{\beta}\in \mathbb{F}^{m\times n}$ by taking the $i$-th column to be $[T(b_i)]_\gamma$:
\[
[T]^{\gamma}_{\beta}\coloneqq
\begin{bmatrix}
\big| &        & \big|\\
[T(b_1)]_\gamma & \cdots & [T(b_n)]_\gamma\\
\big| &        & \big|
\end{bmatrix}.
\]

\end{definition}

The matrix $[T]^{\gamma}_{\beta}$ encodes $T$ completely (relative to the chosen bases): the $j$-th column tells you what $T$ does to the $j$-th basis vector of the domain, expressed in the codomain basis.
Changing bases will change the matrix, but the underlying linear map stays the same.
This is why matrix representations depend on \emph{two} ordered bases: one for the domain and one for the codomain.

\begin{commonmistake}
\begin{itemize}
\item Confusing $[T]^{\gamma}_{\beta}$ (a matrix in $\mathbb{F}^{m\times n}$) with $T$ (an abstract linear map). The matrix is a \emph{representation} that depends on the basis choice.
\item Forgetting to use the correct basis for each side: the input vector must be in $\beta$-coordinates, and the output is in $\gamma$-coordinates.
\item Dimension mismatch: $[T]^{\gamma}_{\beta}$ is $\dim(W)\times\dim(V)$, not $\dim(V)\times\dim(W)$.
\end{itemize}
\end{commonmistake}

\begin{theorem}{Fundamental coordinate identity}{fundamental-coordinate-identity}
With the notation of Definition~\ref{def:matrix-rep}, for every $v\in V$,
\[
[T(v)]_\gamma = [T]^{\gamma}_{\beta}\,[v]_\beta.
\]
\end{theorem}

\begin{proof}
Write $v=\sum_{i=1}^n \lambda_i b_i$, so $[v]_\beta=(\lambda_1,\dots,\lambda_n)^{\mathsf T}$.
By linearity, $T(v)=\sum_{i=1}^n \lambda_i T(b_i)$, hence
\[
[T(v)]_\gamma=\sum_{i=1}^n \lambda_i [T(b_i)]_\gamma
=
\begin{bmatrix}
\big| &        & \big|\\
[T(b_1)]_\gamma & \cdots & [T(b_n)]_\gamma\\
\big| &        & \big|
\end{bmatrix}
\begin{bmatrix}\lambda_1\\ \vdots\\ \lambda_n\end{bmatrix}
=[T]^{\gamma}_{\beta}\,[v]_\beta.
\]
\end{proof}

\begin{theorem}{Composition corresponds to matrix multiplication}{composition-matrix-multiplication}
Let $U,V,W$ be finite-dimensional over $\mathbb{F}$ with ordered bases
$\alpha$ for $U$, $\beta$ for $V$, and $\gamma$ for $W$.
Let $S\in\mathcal{L}(U,V)$ and $T\in\mathcal{L}(V,W)$. Then
\[
[T\circ S]^{\gamma}_{\alpha} = [T]^{\gamma}_{\beta}\,[S]^{\beta}_{\alpha}.
\]
\end{theorem}

\begin{proof}
For all $u\in U$,
\[
[(T\circ S)(u)]_\gamma=[T(S(u))]_\gamma=[T]^{\gamma}_{\beta}[S(u)]_\beta
=[T]^{\gamma}_{\beta}[S]^{\beta}_{\alpha}[u]_\alpha.
\]
Also $[(T\circ S)(u)]_\gamma=[T\circ S]^{\gamma}_{\alpha}[u]_\alpha$.
Since this holds for all $u$, the matrices must be equal.

\end{proof}

The composition formula $[T\circ S]^{\gamma}_{\alpha}=[T]^{\gamma}_{\beta}[S]^{\beta}_{\alpha}$ requires the ``inner'' basis $\beta$ to match: $S$ outputs $\beta$-coordinates, and $T$ reads $\beta$-coordinates. If the intermediate bases do not match, the formula fails.

\begin{examplebox}
\textbf{Worked Example (Tutorial 5, Question 5 style): operator on $\mathcal{P}_2(\mathbb{R})$.}
Define $T\in\mathcal{L}(\mathcal{P}_2(\mathbb{R}))$ by $T(P)=P(X+1)-(P(X^2))''$.
Write $P=a+bX+cX^2$.
\begin{itemize}
\item $P(X+1)=(a+b+c)+(b+2c)X+cX^2$.
\item $(P(X^2))''=(a+bX^2+cX^4)''=2b+12cX^2$.
\end{itemize}
Thus $T(P)=(a-b+c)+(b+2c)X-11cX^2$, giving the standard matrix
\[
[T]_{(1,X,X^2)}=\begin{bmatrix}1&-1&1\\0&1&2\\0&0&-11\end{bmatrix}.
\]
This is upper triangular with diagonal entries $1,1,-11$ (all nonzero), so $\det=-11\neq 0$ and $T$ is bijective.
\end{examplebox}

%====================================================================

\subsection{Change of coordinates and change of matrix representations (Week 7)}

\begin{definition}{Change-of-coordinate matrix}{change-of-coordinates}
Let $V$ be a finite-dimensional vector space over $\mathbb{F}$.
Let $\beta$ and $\gamma$ be ordered bases for $V$.
The \textbf{change-of-coordinate matrix from $\beta$ to $\gamma$} is
\[
[I_V]^{\gamma}_{\beta}\ \coloneqq\ [I_V]^{\gamma}_{\beta},
\]
i.e.\ the matrix representation of the identity map $I_V:V\to V$ with domain basis $\beta$ and codomain basis $\gamma$.

\end{definition}

Why is this natural? Suppose you have a vector $v\in V$ whose $\beta$-coordinates you know, and you want its $\gamma$-coordinates.
The identity map $I_V$ does not change the vector itself, but different bases ``see'' the same vector differently.
The matrix $[I_V]^{\gamma}_{\beta}$ is the translator: it converts $\beta$-coordinates to $\gamma$-coordinates.
Its inverse $[I_V]^{\beta}_{\gamma}=([I_V]^{\gamma}_{\beta})^{-1}$ translates in the opposite direction.

\begin{examtip}
\textbf{Practical computation of $[I_V]^{\gamma}_{\beta}$ (Week 7).}
If $V$ has a convenient ``standard'' basis $\sigma$, form the matrices $P_\beta=[I_V]^{\sigma}_{\beta}$ and $P_\gamma=[I_V]^{\sigma}_{\gamma}$ by stacking basis vectors as columns, then
\[
[I_V]^{\gamma}_{\beta}=P_\gamma^{-1}P_\beta.
\]
An efficient row-reduction approach: form $[\,P_\gamma\mid P_\beta\,]$ and row-reduce to $[\,I_n\mid P_\gamma^{-1}P_\beta\,]$.
\end{examtip}

\begin{theorem}{Changing coordinate vectors}{thm:change-coordinate-vectors}
Let $V$ be finite-dimensional with ordered bases $\beta$ and $\gamma$.
Then for every $v\in V$,
\[
[v]_\gamma = [I_V]^{\gamma}_{\beta}\,[v]_\beta.
\]
\end{theorem}

\begin{proof}
Apply Theorem~\ref{thm:fundamental-coordinate-identity} to $T=I_V$:
\[
[I_V(v)]_\gamma = [I_V]^{\gamma}_{\beta}\,[v]_\beta.
\]
Since $I_V(v)=v$, the result follows.
\end{proof}

\begin{theorem}{Inverse map and inverse matrix (Week 7)}{thm:inverse-matrix-representation}
Let $V,W$ be finite-dimensional over $\mathbb{F}$ with ordered bases $\beta$ for $V$ and $\gamma$ for $W$.
Let $T\in\mathcal{L}(V,W)$. Then $T$ is invertible from $V$ to $W$ if and only if $[T]^{\gamma}_{\beta}$ is an invertible square matrix,
in which case
\[
[T^{-1}]^{\beta}_{\gamma}=\bigl([T]^{\gamma}_{\beta}\bigr)^{-1}.
\]
\end{theorem}

\begin{proof}
($\Rightarrow$) If $T$ is invertible, then $\dim(V)=\dim(W)=n$.
Using composition--multiplication (Theorem~\ref{thm:composition-matrix-multiplication}),
\[
[I_V]^{\beta}_{\beta}=[T^{-1}\circ T]^{\beta}_{\beta}=[T^{-1}]^{\beta}_{\gamma}[T]^{\gamma}_{\beta},
\quad
[I_W]^{\gamma}_{\gamma}=[T\circ T^{-1}]^{\gamma}_{\gamma}=[T]^{\gamma}_{\beta}[T^{-1}]^{\beta}_{\gamma}.
\]
Thus $[T]^{\gamma}_{\beta}$ is invertible with inverse $[T^{-1}]^{\beta}_{\gamma}$.

($\Leftarrow$) Suppose $[T]^{\gamma}_{\beta}\in\mathbb{F}^{n\times n}$ is invertible.
Let $A\coloneqq([T]^{\gamma}_{\beta})^{-1}$.
By the standard existence-and-uniqueness result for a linear map specified on a basis (Weeks 5--6),
there exists $S\in\mathcal{L}(W,V)$ with $[S]^{\beta}_{\gamma}=A$.
Then
\[
[S\circ T]^{\beta}_{\beta}=[S]^{\beta}_{\gamma}[T]^{\gamma}_{\beta}=A[T]^{\gamma}_{\beta}=I_n=[I_V]^{\beta}_{\beta},
\]
and similarly $[T\circ S]^{\gamma}_{\gamma}=I_n=[I_W]^{\gamma}_{\gamma}$.
Hence $S\circ T=I_V$ and $T\circ S=I_W$, so $T$ is invertible and $S=T^{-1}$.
\end{proof}

\begin{theorem}{Change of matrix representations under a change of coordinates (Week 7)}{change-of-matrix-rep}
Let $V$ be finite-dimensional and let $\beta,\gamma$ be ordered bases for $V$.
Let $T\in\mathcal{L}(V)$. Then
\[
[T]^{\gamma}_{\gamma} = [I_V]^{\gamma}_{\beta}\,[T]^{\beta}_{\beta}\,[I_V]^{\beta}_{\gamma}.
\]
\end{theorem}

\begin{proof}
For any $v\in V$,
\[
[T(v)]_\gamma = [I_V]^{\gamma}_{\beta}[T(v)]_\beta
\quad\text{and}\quad
[T(v)]_\beta = [T]^{\beta}_{\beta}[v]_\beta
\quad\text{and}\quad
[v]_\beta = [I_V]^{\beta}_{\gamma}[v]_\gamma.
\]
Substitute the last two into the first:
\[
[T(v)]_\gamma
=[I_V]^{\gamma}_{\beta}[T]^{\beta}_{\beta}[I_V]^{\beta}_{\gamma}[v]_\gamma.
\]
But also $[T(v)]_\gamma=[T]^{\gamma}_{\gamma}[v]_\gamma$.
Since this holds for all $v$, the matrices are equal.

\end{proof}

This formula says: to change the basis from $\beta$ to $\gamma$, you ``sandwich'' the $\beta$-representation between the change-of-coordinate matrices.
Read right-to-left: first convert $\gamma$-coordinates to $\beta$-coordinates (multiply by $[I_V]^{\beta}_{\gamma}$), then apply $T$ in $\beta$-coordinates, then convert back to $\gamma$-coordinates (multiply by $[I_V]^{\gamma}_{\beta}$).
This is exactly the similarity relation $[T]^{\gamma}_{\gamma}=Q^{-1}[T]^{\beta}_{\beta}Q$ with $Q=[I_V]^{\beta}_{\gamma}$.

\begin{keyformula}
\textbf{Computing $[I_V]^{\gamma}_{\beta}$ via a standard basis $\sigma$ (Week 7).}
Let $\sigma$ be a standard ordered basis for $V$ (e.g.\ the standard basis in $\mathbb{F}^n$).
Then
\[
[I_V]^{\gamma}_{\beta} = \bigl([I_V]^{\sigma}_{\gamma}\bigr)^{-1}[I_V]^{\sigma}_{\beta},
\qquad
[I_V]^{\beta}_{\gamma} = \bigl([I_V]^{\sigma}_{\beta}\bigr)^{-1}[I_V]^{\sigma}_{\gamma},
\]
where $[I_V]^{\sigma}_{\beta}=\bigl[\, [b_1]_\sigma\ \cdots\ [b_n]_\sigma\,\bigr]$ is formed by stacking basis vectors (in $\sigma$-coordinates) as columns.
\end{keyformula}

\begin{examplebox}
\textbf{Worked Example (Lecture Week 7 problem): compute $[I_V]^{\alpha}_{\beta}$ and coordinates.}

Let $\alpha=(a_1,a_2,a_3)$ and $\beta=(b_1,b_2,b_3)$ be ordered bases of $V$ such that
\[
b_1=2a_1-a_2+a_3,\qquad b_2=3a_2+a_3,\qquad b_3=-3a_1+2a_3.
\]
\textbf{(1) Compute $[I_V]^{\alpha}_{\beta}$.}
By definition, the $i$-th column of $[I_V]^{\alpha}_{\beta}$ is $[b_i]_\alpha$.
Thus
\[
[b_1]_\alpha=\begin{bmatrix}2\\-1\\1\end{bmatrix},\quad
[b_2]_\alpha=\begin{bmatrix}0\\3\\1\end{bmatrix},\quad
[b_3]_\alpha=\begin{bmatrix}-3\\0\\2\end{bmatrix},
\]
so
\[
[I_V]^{\alpha}_{\beta}=
\begin{bmatrix}
2 & 0 & -3\\
-1 & 3 & 0\\
1 & 1 & 2
\end{bmatrix}.
\]
\textbf{(2) Compute $[b_1-2b_2+2b_3]_\alpha$.}
Using linearity of coordinate expansion in the basis $\alpha$,
\[
[b_1-2b_2+2b_3]_\alpha=[b_1]_\alpha-2[b_2]_\alpha+2[b_3]_\alpha
=
\begin{bmatrix}2\\-1\\1\end{bmatrix}
-2\begin{bmatrix}0\\3\\1\end{bmatrix}
+2\begin{bmatrix}-3\\0\\2\end{bmatrix}
=
\begin{bmatrix}-4\\-7\\3\end{bmatrix}.
\]
\end{examplebox}

\begin{examplebox}
\textbf{Worked Example (Lecture Week 7 problem): compute change-of-coordinate matrices in $\mathbb{R}^2$.}

Let $\beta=(({-9},1),({-5},{-1}))$ and $\gamma=((1,{-4}),(3,{-5}))$ be ordered bases for $\mathbb{R}^2$.
Let $\sigma$ be the standard basis. Then
\[
[I_{\mathbb{R}^2}]^{\sigma}_{\beta}
=
\begin{bmatrix}
-9 & -5\\
1 & -1
\end{bmatrix},
\qquad
[I_{\mathbb{R}^2}]^{\sigma}_{\gamma}
=
\begin{bmatrix}
1 & 3\\
-4 & -5
\end{bmatrix}.
\]
Hence
\[
[I_{\mathbb{R}^2}]^{\gamma}_{\beta}
=\bigl([I_{\mathbb{R}^2}]^{\sigma}_{\gamma}\bigr)^{-1}[I_{\mathbb{R}^2}]^{\sigma}_{\beta}
=
\begin{bmatrix}
6 & 4\\
-5 & -3
\end{bmatrix},
\qquad
[I_{\mathbb{R}^2}]^{\beta}_{\gamma}
=\bigl([I_{\mathbb{R}^2}]^{\sigma}_{\beta}\bigr)^{-1}[I_{\mathbb{R}^2}]^{\sigma}_{\gamma}
=
\begin{bmatrix}
-\tfrac{3}{2} & -2\\
\tfrac{5}{2} & 3
\end{bmatrix}.
\]
(Indeed $[I]^{\gamma}_{\beta}[I]^{\beta}_{\gamma}=I_2$.)
\end{examplebox}

\begin{examplebox}
\textbf{Worked Example (Tutorial 6, Question 7): compute $[T]^{\gamma}_{\gamma}$ by similarity.}
Define $T\in\mathcal{L}(\mathbb{R}^3)$ by $T(v)=Av$ where
$A=\begin{bmatrix}2&1&0\\1&1&3\\0&-1&0\end{bmatrix}$.
Let $\gamma=(g_1,g_2,g_3)$ with
$g_1=\begin{bmatrix}-1\\0\\0\end{bmatrix}$, $g_2=\begin{bmatrix}2\\1\\0\end{bmatrix}$, $g_3=\begin{bmatrix}1\\1\\1\end{bmatrix}$.

Form $P=[g_1\; g_2\; g_3]=\begin{bmatrix}-1&2&1\\0&1&1\\0&0&1\end{bmatrix}$. Then
\[
[T]^{\gamma}_{\gamma}=P^{-1}AP=\begin{bmatrix}0&2&8\\-1&4&6\\0&-1&-1\end{bmatrix}.
\]
\textbf{Sanity check:} $\operatorname{tr}([T]^{\gamma}_{\gamma})=0+4+(-1)=3=\operatorname{tr}(A)=2+1+0$ \checkmark, and $\det([T]^{\gamma}_{\gamma})=\det(A)$ (both can be verified numerically).
\end{examplebox}

\begin{examplebox}
\textbf{Worked Example (Lecture Week 7/8 problem): diagonalizing a map via a tailored basis.}

Define $T\in\mathcal{L}(\mathbb{R}^3)$ by
\[
T(x,y,z)=(2x+y,\ 5y+3z,\ 8z).
\]
Let $\beta=(b_1,b_2,b_3)$ with $b_1=(1,0,0)$, $b_2=(1,3,0)$, $b_3=(1,6,6)$.
\begin{enumerate}
\item Compute $[T]^{\beta}_{\beta}$.
\item Compute $T^{4}(0,-3,6)$.
\end{enumerate}

\medskip
\textbf{Solution.}
\textbf{(1)} Compute $T(b_i)$ and express in the basis $\beta$:
\[
T(b_1)=(2,0,0)=2b_1,\qquad
T(b_2)=(5,15,0)=5b_2,\qquad
T(b_3)=(8,48,48)=8b_3.
\]
Thus
\[
[T]^{\beta}_{\beta}=
\begin{bmatrix}
2&0&0\\
0&5&0\\
0&0&8
\end{bmatrix}.
\]
\textbf{(2)} First compute $[(0,-3,6)]_\beta$. Solve
\[
(0,-3,6)=\alpha b_1+\beta b_2+\gamma b_3=(\alpha+\beta+\gamma,\ 3\beta+6\gamma,\ 6\gamma),
\]
giving $\gamma=1$, $\beta=-3$, $\alpha=2$. Hence $[(0,-3,6)]_\beta=(2,-3,1)^{\mathsf T}$.
Then
\[
[T^4(0,-3,6)]_\beta=\bigl([T]^{\beta}_{\beta}\bigr)^4[(0,-3,6)]_\beta
=\mathrm{diag}(2^4,5^4,8^4)\begin{bmatrix}2\\-3\\1\end{bmatrix}
=\begin{bmatrix}32\\-1875\\4096\end{bmatrix}.
\]
Convert back:
\[
T^4(0,-3,6)=32b_1-1875b_2+4096b_3=(2253,\ 18951,\ 24576).
\]
\end{examplebox}

\begin{extension}[title={Extension (Beyond Syllabus: the coordinate isomorphism $[\,\cdot\,]_\beta$ packages everything)}]
\textbf{Syllabus link:} MH1201 Weeks 7--8 — Change of coordinates.

\textbf{Theorem.} Let $U,V$ be finite-dimensional over $\mathbb{F}$ with ordered bases $\alpha$ (for $U$) and $\beta$ (for $V$).
Define the map
\[
\Phi_{\beta,\alpha}:\mathcal{L}(U,V)\to \mathbb{F}^{\dim(V)\times \dim(U)},\qquad
\Phi_{\beta,\alpha}(S)=[S]^{\beta}_{\alpha}.
\]
Then $\Phi_{\beta,\alpha}$ is an isomorphism of vector spaces. Moreover, for $S\in\mathcal{L}(U,V)$ and $T\in\mathcal{L}(V,W)$
(with ordered bases $\gamma$ for $W$),
\[
[T\circ S]^{\gamma}_{\alpha}=[T]^{\gamma}_{\beta}[S]^{\beta}_{\alpha}.
\]

\textbf{Proof.}
\emph{Linearity:} For any $S_1,S_2\in\mathcal{L}(U,V)$ and $\lambda\in\mathbb{F}$ and any $u\in U$,
\[
[(S_1+S_2)(u)]_\beta=[S_1(u)+S_2(u)]_\beta=[S_1(u)]_\beta+[S_2(u)]_\beta,
\]
so by the fundamental identity (Theorem~\ref{thm:fundamental-coordinate-identity}) applied to each map,
\[
[S_1+S_2]^{\beta}_{\alpha}[u]_\alpha=[S_1]^{\beta}_{\alpha}[u]_\alpha+[S_2]^{\beta}_{\alpha}[u]_\alpha.
\]
Since this holds for all $u$, $[S_1+S_2]^{\beta}_{\alpha}=[S_1]^{\beta}_{\alpha}+[S_2]^{\beta}_{\alpha}$.
Similarly $[\lambda S_1]^{\beta}_{\alpha}=\lambda [S_1]^{\beta}_{\alpha}$.

\emph{Injective:} If $[S]^{\beta}_{\alpha}=0$, then for all $u$, $[S(u)]_\beta=0$, hence $S(u)=0$. So $S=0$.

\emph{Surjective:} Given any matrix $A\in\mathbb{F}^{\dim(V)\times \dim(U)}$, define $S$ on the basis $\alpha=(a_1,\dots,a_n)$ by
requiring $[S(a_j)]_\beta$ to be the $j$-th column of $A$, then extend linearly.
This produces a linear map with matrix $[S]^{\beta}_{\alpha}=A$ (uniqueness/existence from basis specification).

Finally, the composition identity is exactly Theorem~\ref{thm:composition-matrix-multiplication}. \qed
\end{extension}

\begin{commonmistake}
\begin{itemize}
\item Swapping the direction of $[I_V]^{\gamma}_{\beta}$: it converts \emph{$\beta$-coordinates to $\gamma$-coordinates} via $[v]_\gamma=[I_V]^{\gamma}_{\beta}[v]_\beta$.
\item Forgetting that $[I_V]^{\beta}_{\gamma}=\bigl([I_V]^{\gamma}_{\beta}\bigr)^{-1}$ (invertibility holds because $I_V$ is invertible).
\end{itemize}
\end{commonmistake}

%====================================================================

\subsection{Similarity of square matrices (Week 8)}

\begin{definition}{Similarity}{similarity}
Let $\mathbb{F}$ be a field and $n\in\mathbb{N}$.
For $A,B\in\mathbb{F}^{n\times n}$, we say \textbf{$A$ is similar to $B$}, written $A\sim B$, if there exists an invertible
$Q\in\mathbb{F}^{n\times n}$ such that
\[
B=Q^{-1}AQ.
\]
\end{definition}

\begin{theorem}{Similarity is an equivalence relation}{similarity-equivalence}
Let $\mathbb{F}$ be a field and $n\in\mathbb{N}$. The relation $\sim$ on $\mathbb{F}^{n\times n}$ is an equivalence relation.
\end{theorem}

\begin{proof}
\textbf{Reflexive.} For any $A$, take $Q=I_n$, then $A=I_n^{-1}AI_n$, so $A\sim A$.

\textbf{Symmetric.} If $A\sim B$, then $B=Q^{-1}AQ$ for some invertible $Q$.
Multiply by $Q$ on the left and $Q^{-1}$ on the right to get $A=QBQ^{-1}$, so $B\sim A$.

\textbf{Transitive.} If $A\sim B$ and $B\sim C$, then $B=Q^{-1}AQ$ and $C=P^{-1}BP$ for invertible $P,Q$.
Substitute:
\[
C=P^{-1}(Q^{-1}AQ)P=(QP)^{-1}A(QP),
\]
so $A\sim C$.

\end{proof}

Similarity partitions $\mathbb{F}^{n\times n}$ into equivalence classes.
Each class consists of all matrices that represent the \emph{same} linear operator under various basis choices.
This is why similarity invariants (determinant, trace, eigenvalues, characteristic polynomial) are really properties of the underlying operator $T$, not of any particular matrix.

\begin{theorem}{Determinant and trace are similarity invariants}{similarity-invariants}
Let $\mathbb{F}$ be a field and $A,B\in\mathbb{F}^{n\times n}$ with $A\sim B$. Then
\[
\det(A)=\det(B)\qquad\text{and}\qquad \operatorname{tr}(A)=\operatorname{tr}(B).
\]
\end{theorem}

\begin{proof}
Let $B=Q^{-1}AQ$ with $Q$ invertible.
Using multiplicativity of determinant,
\[
\det(B)=\det(Q^{-1}AQ)=\det(Q^{-1})\det(A)\det(Q)=\det(Q)^{-1}\det(A)\det(Q)=\det(A).
\]
For trace, use $\operatorname{tr}(XY)=\operatorname{tr}(YX)$ (valid for square matrices of the same size):
\[
\operatorname{tr}(B)=\operatorname{tr}(Q^{-1}AQ)=\operatorname{tr}(AQQ^{-1})=\operatorname{tr}(A).
\]
\end{proof}

\begin{theorem}{Matrix representations under different ordered bases are similar}{matrix-rep-similar}
Let $V$ be finite-dimensional over $\mathbb{F}$, and let $\beta,\gamma$ be ordered bases for $V$.
Let $T\in\mathcal{L}(V)$. Then
\[
[T]^{\beta}_{\beta}\sim [T]^{\gamma}_{\gamma}.
\]
\end{theorem}

\begin{proof}
By Theorem~\ref{thm:change-of-matrix-rep},
\[
[T]^{\gamma}_{\gamma}=[I_V]^{\gamma}_{\beta}[T]^{\beta}_{\beta}[I_V]^{\beta}_{\gamma}.
\]
Since $[I_V]^{\beta}_{\gamma}=\bigl([I_V]^{\gamma}_{\beta}\bigr)^{-1}$ (because $I_V$ is invertible),
this is exactly the form $Q^{-1}AQ$ with $Q=[I_V]^{\beta}_{\gamma}$ (equivalently $Q^{-1}=[I_V]^{\gamma}_{\beta}$).
Thus the two matrices are similar.
\end{proof}

\begin{commonmistake}
\begin{itemize}
\item Confusing similarity with equality: $A\sim B$ means they represent the \emph{same linear operator} in different coordinates (bases), not that $A=B$.
\item When proving invariants, do not assume $\operatorname{tr}(XY)=\operatorname{tr}(X)\operatorname{tr}(Y)$ (false). The correct identity is $\operatorname{tr}(XY)=\operatorname{tr}(YX)$.
\end{itemize}
\end{commonmistake}

%====================================================================

\subsection{Eigenvalues, eigenvectors, eigenspaces (Week 8)}

\begin{definition}{Eigenvalue}{def:eigenvalue}
Let $V$ be a vector space over a field $\mathbb{F}$ and let $T\in\mathcal{L}(V)$.
A scalar $\lambda\in\mathbb{F}$ is an \textbf{eigenvalue} of $T$ if any (hence all) of the following equivalent statements hold:
\begin{enumerate}
\item $\operatorname{Null}(T-\lambda I_V)\neq \{0\}$;
\item $\{0\}\subsetneq \operatorname{Null}(T-\lambda I_V)$;
\item $T-\lambda I_V$ is not injective.
\end{enumerate}

\end{definition}

Eigenvalues are ``field-specific'': they must come from the base field $\mathbb{F}$ of the vector space.
The same matrix can have eigenvalues over $\mathbb{C}$ but none over $\mathbb{R}$ (see the rotation example below).
This is why the problem statement should always specify the field.

Intuitively, $\lambda$ is an eigenvalue if $T-\lambda I_V$ ``collapses'' some nonzero vector to zero, i.e.\ $T$ acts on some direction simply as scaling by $\lambda$.

\begin{keyformula}
\textbf{Field-specific warning (Week 8).}
Eigenvalues must lie in the base field $\mathbb{F}$ of the vector space.
The same linear map can have no eigenvalues over $\mathbb{R}$ but have eigenvalues over $\mathbb{C}$.
\end{keyformula}

\begin{theorem}{Zero eigenvalue $\Longleftrightarrow$ non-injectivity}{zero-eig-noninj}
Let $V$ be a vector space over $\mathbb{F}$ and $T\in\mathcal{L}(V)$.
Then $0\in\mathbb{F}$ is an eigenvalue of $T$ if and only if $T$ is not injective.
\end{theorem}

\begin{proof}
$0$ is an eigenvalue $\Longleftrightarrow \operatorname{Null}(T-0\cdot I_V)=\operatorname{Null}(T)\neq\{0\}
\Longleftrightarrow T$ is not injective by Theorem~\ref{thm:null-space-test-injective}.
\end{proof}

\begin{definition}{Eigenvector and eigenspace}{def:eigenvector-eigenspace}
Let $V$ be a vector space over $\mathbb{F}$ and let $T\in\mathcal{L}(V)$.
A nonzero vector $v\in V$ is an \textbf{eigenvector} of $T$ if there exists $\lambda\in\mathbb{F}$ such that
\[
T(v)=\lambda v.
\]
For an eigenvalue $\lambda$ of $T$, the \textbf{eigenspace} of $T$ for $\lambda$ is
\[
E_{T,\lambda}\coloneqq \operatorname{Null}(T-\lambda I_V).
\]
Then the eigenvectors for $\lambda$ are exactly $E_{T,\lambda}\setminus\{0\}$.

\end{definition}

The eigenspace $E_{T,\lambda}$ is a subspace of $V$ (it is the null space of $T-\lambda I_V$).
Its nonzero elements are the eigenvectors for $\lambda$.
The zero vector is in $E_{T,\lambda}$ but is \emph{never} called an eigenvector (by convention, eigenvectors must be nonzero).
Note that the eigenvalue associated to a given eigenvector $v$ is unique: if $T(v)=\lambda v$ and $T(v)=\mu v$ with $v\neq 0$, then $(\lambda-\mu)v=0$ implies $\lambda=\mu$.

\begin{theorem}{Eigenvector--kernel equivalence}{thm:eig-kernel}
Let $T\in\mathcal{L}(V)$ and $\lambda\in\mathbb{F}$. For $v\in V$,
\[
v\in E_{T,\lambda}\ \Longleftrightarrow\ (T-\lambda I_V)(v)=0\ \Longleftrightarrow\ T(v)=\lambda v.
\]
\end{theorem}

\begin{proof}
Immediate from expanding $(T-\lambda I_V)(v)=T(v)-\lambda I_V(v)=T(v)-\lambda v$.
\end{proof}

\begin{theorem}{Computing eigenvalues via a matrix representation (finite-dimensional)}{thm:eig-det}
Let $V$ be finite-dimensional over $\mathbb{F}$ with $\dim(V)=n$, and let $T\in\mathcal{L}(V)$.
Let $\beta$ be any ordered basis for $V$.
Then $\lambda\in\mathbb{F}$ is an eigenvalue of $T$ if and only if
\[
\det\bigl([T]^{\beta}_{\beta}-\lambda I_n\bigr)=0.
\]
\end{theorem}

\begin{proof}
$\lambda$ is an eigenvalue $\Longleftrightarrow T-\lambda I_V$ is not injective
$\Longleftrightarrow [T-\lambda I_V]^{\beta}_{\beta}$ is not invertible
$\Longleftrightarrow [T]^{\beta}_{\beta}-\lambda[I_V]^{\beta}_{\beta}$ is not invertible
$\Longleftrightarrow [T]^{\beta}_{\beta}-\lambda I_n$ is not invertible
$\Longleftrightarrow \det([T]^{\beta}_{\beta}-\lambda I_n)=0$.
\end{proof}

\begin{definition}{Characteristic polynomial function}{charpoly}
Let $V$ be finite-dimensional over $\mathbb{F}$ with $\dim(V)=n$ and let $T\in\mathcal{L}(V)$.
Fix any ordered basis $\beta$ of $V$.
Define the \textbf{characteristic polynomial function} of $T$ by
\[
p_T(\lambda)\coloneqq \det\bigl([T]^{\beta}_{\beta}-\lambda I_n\bigr)\qquad(\lambda\in\mathbb{F}).
\]

\end{definition}

For an $n\times n$ matrix $A$, the characteristic polynomial $p_T(\lambda)=\det(A-\lambda I_n)$ is a polynomial of degree $n$ in $\lambda$.
\begin{itemize}
\item For $n=2$: $\det\begin{bmatrix}a-\lambda&b\\c&d-\lambda\end{bmatrix}=\lambda^2-(a+d)\lambda+(ad-bc)=\lambda^2-\operatorname{tr}(A)\lambda+\det(A)$.
\item For $n=3$: expand using cofactor expansion along a row or column.
\end{itemize}
The roots of $p_T$ in $\mathbb{F}$ are the eigenvalues. If $\mathbb{F}=\mathbb{R}$ and the polynomial has no real roots, then $T$ has no eigenvalues over $\mathbb{R}$.

\begin{theorem}{Basis-independence of $p_T$}{charpoly-basis-indep}
Let $V$ be finite-dimensional and $T\in\mathcal{L}(V)$.
If $\beta$ and $\gamma$ are two ordered bases of $V$, then for all $\lambda\in\mathbb{F}$,
\[
\det\bigl([T]^{\beta}_{\beta}-\lambda I_n\bigr)=\det\bigl([T]^{\gamma}_{\gamma}-\lambda I_n\bigr).
\]
Hence $p_T$ is well-defined (independent of the chosen ordered basis).
\end{theorem}

\begin{proof}
By Theorem~\ref{thm:matrix-rep-similar}, there exists invertible $Q$ such that
$[T]^{\gamma}_{\gamma}=Q^{-1}[T]^{\beta}_{\beta}Q$.
Then for every $\lambda$,
\[
[T]^{\gamma}_{\gamma}-\lambda I_n = Q^{-1}\bigl([T]^{\beta}_{\beta}-\lambda I_n\bigr)Q,
\]
so the two matrices are similar. By Theorem~\ref{thm:similarity-invariants}, similar matrices have equal determinants.
\end{proof}

\begin{keyformula}
\textbf{Eigenvalues = roots of $p_T$ (finite-dimensional).}
If $\dim(V)=n$, then eigenvalues of $T$ are exactly the $\mathbb{F}$-roots of $p_T$.
In particular, $T$ has at most $n$ distinct eigenvalues in $\mathbb{F}$.
\end{keyformula}

\begin{examplebox}
\textbf{Worked Example (Lecture Week 8): eigenvalues of $0_{V\to V}$ and $I_V$.}
Let $V$ be nonzero over $\mathbb{F}$.
\begin{itemize}
\item For the zero map $0_{V\to V}$, we have $0_{V\to V}(v)=0$ for all $v$, so every nonzero $v$ satisfies $0(v)=0\cdot v$.
Thus the only eigenvalue is $\lambda=0$.
\item For the identity $I_V$, we have $I_V(v)=1\cdot v$ for all $v\neq 0$, so the only eigenvalue is $\lambda=1$.
\end{itemize}
\end{examplebox}

\begin{examplebox}
\textbf{Worked Example (Lecture Week 8): rotation matrix over $\mathbb{R}$ vs $\mathbb{C}$.}
Let $A=\begin{bmatrix}0&1\\-1&0\end{bmatrix}$ and define $T(v)=Av$.
Then
\[
\det(A-\lambda I_2)=\det\begin{bmatrix}-\lambda&1\\-1&-\lambda\end{bmatrix}=\lambda^2+1.
\]
Over $\mathbb{R}$, $\lambda^2+1=0$ has no roots, so $T$ has no real eigenvalues.
Over $\mathbb{C}$, the roots are $\lambda=\pm i$, so $T$ has eigenvalues $\pm i$ in $\mathbb{C}$.
\end{examplebox}

\begin{examplebox}
\textbf{Worked Example (Lecture Week 8): verify an eigenvector quickly.}
Let $T(v)=Av$ on $\mathbb{R}^3$ where
\[
A=\begin{bmatrix}
0&6&8\\[2pt]
\frac12&0&0\\[2pt]
0&\frac12&0
\end{bmatrix}.
\]
Check $v=\begin{bmatrix}16\\4\\1\end{bmatrix}$:
\[
Av=
\begin{bmatrix}
0\cdot 16+6\cdot 4+8\cdot 1\\
\frac12\cdot 16\\
\frac12\cdot 4
\end{bmatrix}
=
\begin{bmatrix}32\\8\\2\end{bmatrix}
=2\begin{bmatrix}16\\4\\1\end{bmatrix}.
\]
Hence $v$ is an eigenvector with eigenvalue $\lambda=2$.
\end{examplebox}

\begin{examplebox}
\textbf{Worked Example (Lecture Week 8): two eigenvectors without heavy computation.}
Let $T(v)=Av$ on $\mathbb{R}^3$ where
\[
A=\begin{bmatrix}
1&2&3\\
1&2&3\\
1&2&3
\end{bmatrix}.
\]
Let $u=(1,1,1)^{\mathsf T}$. Then
\[
Au=\begin{bmatrix}6\\6\\6\end{bmatrix}=6u,
\]
so $u$ is an eigenvector for $\lambda=6$.
Now take any nonzero $w$ with $(1,2,3)\cdot w=0$, e.g.\ $w=(2,-1,0)^{\mathsf T}$.
Then $Aw=u\,(1,2,3)\cdot w=0$, so $w$ is an eigenvector for $\lambda=0$.
Thus we found two eigenvectors for distinct eigenvalues $6$ and $0$.
\end{examplebox}

\begin{theorem}{Eigenvectors for distinct eigenvalues are linearly independent}{distinct-eigs-LI}
Let $V$ be a vector space over $\mathbb{F}$ and let $T\in\mathcal{L}(V)$.
Let $\lambda_1,\dots,\lambda_n$ be distinct eigenvalues of $T$ and let $v_i\neq 0$ be an eigenvector for $\lambda_i$.
Then $\{v_1,\dots,v_n\}$ is linearly independent.
\end{theorem}

\begin{proof}
We use induction on $n$.

For $n=1$, $\{v_1\}$ is independent since $v_1\neq 0$.

Assume the statement true for $n-1$ and consider a relation
\[
\mu_1 v_1+\cdots+\mu_n v_n=0.
\]
Apply $T$:
\[
0=T(0)=\mu_1 T(v_1)+\cdots+\mu_n T(v_n)=\mu_1\lambda_1 v_1+\cdots+\mu_n\lambda_n v_n.
\]
Multiply the original relation by $\lambda_n$ and subtract:
\[
0=\mu_1(\lambda_1-\lambda_n)v_1+\cdots+\mu_{n-1}(\lambda_{n-1}-\lambda_n)v_{n-1}.
\]
By distinctness, each $(\lambda_i-\lambda_n)\neq 0$, so by the induction hypothesis $\{v_1,\dots,v_{n-1}\}$ is independent and hence
\[
\mu_1(\lambda_1-\lambda_n)=\cdots=\mu_{n-1}(\lambda_{n-1}-\lambda_n)=0 \ \Longrightarrow\ \mu_1=\cdots=\mu_{n-1}=0.
\]
Then the original relation becomes $\mu_n v_n=0$, so $\mu_n=0$ since $v_n\neq 0$.
Thus all coefficients are zero, proving independence.
\end{proof}

\begin{theorem}{Sum of distinct eigenspaces is a direct sum}{eigenspaces-direct-sum}
Let $V$ be a vector space over $\mathbb{F}$ and let $T\in\mathcal{L}(V)$.
Let $\lambda_1,\dots,\lambda_n$ be distinct eigenvalues. Then
\[
E_{T,\lambda_1}+\cdots+E_{T,\lambda_n}=E_{T,\lambda_1}\oplus\cdots\oplus E_{T,\lambda_n}.
\]
\end{theorem}

\begin{proof}
We show uniqueness of zero decomposition. Let $v_i\in E_{T,\lambda_i}$ and suppose
\[
v_1+\cdots+v_n=0.
\]
If all $v_i=0$ we are done. Otherwise remove the zero terms and keep a nonempty sub-list of nonzero vectors,
still satisfying a linear combination equal to $0$:
\[
v_{i_1}+\cdots+v_{i_k}=0,\qquad v_{i_j}\neq 0.
\]
But each $v_{i_j}$ is an eigenvector for the distinct eigenvalue $\lambda_{i_j}$, so by
Theorem~\ref{thm:distinct-eigs-LI} the set $\{v_{i_1},\dots,v_{i_k}\}$ is linearly independent, contradicting
a nontrivial dependence relation. Hence all $v_i=0$, so the sum is direct.

\end{proof}

This theorem is the foundation of diagonalization: if the eigenspaces span the whole space, then we can choose a basis of eigenvectors.
The direct-sum property ensures that eigenvectors from different eigenspaces never ``interfere'' with each other, so combining bases of individual eigenspaces always yields a linearly independent set.

\begin{examplebox}
\textbf{Worked Example (Lecture Week 8): eigenvalues/eigenspaces of a ``swap'' operator on $\mathbb{R}^{2,2}$.}
Define $T\in\mathcal{L}(\mathbb{R}^{2,2})$ by
\[
T\!\left(\begin{bmatrix}a&b\\c&d\end{bmatrix}\right)=
\begin{bmatrix}c&d\\a&b\end{bmatrix}.
\]
Then $T^2=I$, so any eigenvalue $\lambda$ must satisfy $\lambda^2=1$, hence $\lambda\in\{1,-1\}$.

\textbf{Eigenspace for $\lambda=1$.} Solve $T(A)=A$:
\[
\begin{bmatrix}c&d\\a&b\end{bmatrix}=\begin{bmatrix}a&b\\c&d\end{bmatrix}
\Longleftrightarrow a=c,\ b=d,
\]
so
\[
E_{T,1}=\left\{\begin{bmatrix}a&b\\a&b\end{bmatrix}:a,b\in\mathbb{R}\right\}.
\]
\textbf{Eigenspace for $\lambda=-1$.} Solve $T(A)=-A$:
\[
\begin{bmatrix}c&d\\a&b\end{bmatrix}=\begin{bmatrix}-a&-b\\-c&-d\end{bmatrix}
\Longleftrightarrow c=-a,\ d=-b,
\]
so
\[
E_{T,-1}=\left\{\begin{bmatrix}a&b\\-a&-b\end{bmatrix}:a,b\in\mathbb{R}\right\}.
\]
Moreover, $\mathbb{R}^{2,2}=E_{T,1}\oplus E_{T,-1}$ by Theorem~\ref{thm:eigenspaces-direct-sum}.
\end{examplebox}

\begin{examplebox}
\textbf{Worked Example (Tutorial 7, Question 5): eigenvalues/eigenspaces of $F(x,y)=(y,x+y)$.}
Define $F\in\mathcal{L}(\mathbb{R}^{2,1})$ by
\[
F\!\left(\begin{bmatrix}x\\y\end{bmatrix}\right)=\begin{bmatrix}y\\x+y\end{bmatrix}
=\begin{bmatrix}0&1\\1&1\end{bmatrix}\begin{bmatrix}x\\y\end{bmatrix}.
\]
Let $A=\begin{bmatrix}0&1\\1&1\end{bmatrix}$. Then
\[
\det(A-\lambda I_2)=\det\begin{bmatrix}-\lambda&1\\1&1-\lambda\end{bmatrix}=\lambda^2-\lambda-1,
\]
so eigenvalues are
\[
\lambda_{\pm}=\frac{1\pm\sqrt{5}}{2}.
\]
For $\lambda=\lambda_{\pm}$, solve $(A-\lambda I_2)v=0$:
\[
\begin{bmatrix}-\lambda&1\\1&1-\lambda\end{bmatrix}\begin{bmatrix}x\\y\end{bmatrix}=0
\Longleftrightarrow y=\lambda x.
\]
Hence
\[
E_{F,\lambda}=\operatorname{span}\left\{\begin{bmatrix}1\\ \lambda\end{bmatrix}\right\}
\quad\text{for }\lambda\in\left\{\frac{1+\sqrt5}{2},\frac{1-\sqrt5}{2}\right\}.
\]
\end{examplebox}

\begin{extension}[title={Extension (Tutorial 7, Question 4, Method 2)}]
\textbf{Syllabus link:} MH1201 Weeks 7--8 — Eigenvalues and eigenvectors.

\textbf{Lemma (rotation eigenvalues over $\mathbb{R}$).}
Fix $\theta\in\mathbb{R}$ and define $R_\theta\in\mathcal{L}(\mathbb{R}^2)$ by
\[
R_\theta(x,y)=(x\cos\theta-y\sin\theta,\ x\sin\theta+y\cos\theta).
\]
If $\theta\notin\mathbb{Z}\pi$, then $R_\theta$ has no real eigenvalues.

\textbf{Proof.}
Suppose for contradiction that $\lambda\in\mathbb{R}$ is an eigenvalue with eigenvector $u\neq 0$.
Then $R_\theta(u)=\lambda u$ implies $R_\theta(u)$ is a scalar multiple of $u$, hence $R_\theta(u)$ is parallel to $u$.
But geometrically $R_\theta$ rotates every nonzero vector by angle $\theta$.
A rotation maps a nonzero vector to a parallel vector if and only if the rotation angle is $0$ or $\pi$ modulo $2\pi$,
i.e.\ $\theta\in\mathbb{Z}\pi$. This contradicts $\theta\notin\mathbb{Z}\pi$.
Hence no real eigenvalue exists. \qed
\end{extension}

\begin{extension}[title={Extension (Tutorial 7, Question 6, Method 2)}]
\textbf{Syllabus link:} MH1201 Weeks 7--8 — Eigenvalues and eigenspaces as invariant subspaces.

\textbf{Lemma (invariant decomposition yields eigenstructure).}
Let $T\in\mathcal{L}(\mathbb{R}^{2,2})$ be defined by
\[
T(A)=A^{\mathsf T}+2\,\operatorname{tr}(A)\,I_2.
\]
Define
\[
\mathrm{Sym}_2=\{A:A^{\mathsf T}=A\},\qquad \mathrm{Skew}_2=\{A:A^{\mathsf T}=-A\}.
\]
Then $\mathbb{R}^{2,2}=\mathrm{Sym}_2\oplus \mathrm{Skew}_2$, and:
\[
E_{T,-1}=\mathrm{Skew}_2,\qquad
E_{T,5}=\mathbb{R}I_2,\qquad
E_{T,1}=\{A\in\mathrm{Sym}_2:\operatorname{tr}(A)=0\}.
\]
In particular, $\mathbb{R}^{2,2}$ has a basis of eigenvectors of $T$.

\textbf{Proof.}
\emph{Step 1: direct sum decomposition.}
For any $A$, write
\[
A=\frac12(A+A^{\mathsf T})+\frac12(A-A^{\mathsf T}),
\]
where $\frac12(A+A^{\mathsf T})\in\mathrm{Sym}_2$ and $\frac12(A-A^{\mathsf T})\in\mathrm{Skew}_2$.
Also $\mathrm{Sym}_2\cap \mathrm{Skew}_2=\{0\}$, so $\mathbb{R}^{2,2}=\mathrm{Sym}_2\oplus\mathrm{Skew}_2$.

\emph{Step 2: eigenvalue on $\mathrm{Skew}_2$.}
If $A\in\mathrm{Skew}_2$, then $A^{\mathsf T}=-A$ and $\operatorname{tr}(A)=0$, hence
\[
T(A)=A^{\mathsf T}+2\,\operatorname{tr}(A)I_2=-A.
\]
Thus every nonzero $A\in\mathrm{Skew}_2$ is an eigenvector with eigenvalue $-1$, and $E_{T,-1}=\mathrm{Skew}_2$.

\emph{Step 3: reduce on $\mathrm{Sym}_2$ and split by trace.}
If $A\in\mathrm{Sym}_2$, then $A^{\mathsf T}=A$, so $T(A)=A+2\,\operatorname{tr}(A)I_2$.
Now decompose
\[
\mathrm{Sym}_2=\mathbb{R}I_2\oplus \mathrm{Sym}_2^0,\qquad \mathrm{Sym}_2^0\coloneqq\{A\in\mathrm{Sym}_2:\operatorname{tr}(A)=0\}.
\]
This is a direct sum since $\operatorname{tr}(I_2)=2\neq 0$ and $\operatorname{tr}(\cdot)=0$ on $\mathrm{Sym}_2^0$.
If $A\in\mathrm{Sym}_2^0$, then $T(A)=A$, so $A\in E_{T,1}$ and $E_{T,1}=\mathrm{Sym}_2^0$.
If $A=tI_2$, then $\operatorname{tr}(A)=2t$ and
\[
T(A)=tI_2+2(2t)I_2=5tI_2,
\]
so $\mathbb{R}I_2\subseteq E_{T,5}$ and equality holds by the decomposition (no other symmetric trace-nonzero directions remain).

\emph{Step 4: eigenbasis.}
Choose bases of $\mathrm{Skew}_2$, $\mathbb{R}I_2$, and $\mathrm{Sym}_2^0$; their union is an eigenbasis because the sum is direct and spans $\mathbb{R}^{2,2}$.
For example,
\[
\mathrm{Skew}_2=\operatorname{span}\!\left\{\begin{bmatrix}0&1\\-1&0\end{bmatrix}\right\},\quad
\mathbb{R}I_2=\operatorname{span}\!\left\{\begin{bmatrix}1&0\\0&1\end{bmatrix}\right\},\quad
\mathrm{Sym}_2^0=\operatorname{span}\!\left\{\begin{bmatrix}1&0\\0&-1\end{bmatrix},\begin{bmatrix}0&1\\1&0\end{bmatrix}\right\}.
\]
\qed
\end{extension}

\begin{examtip}
\textbf{Weeks 7--8 exam patterns.}
\begin{itemize}
\item \textbf{Change-of-coordinates:} compute $[I_V]^{\gamma}_{\beta}$ by stacking basis vectors in a standard basis and inverting once.
Then convert coordinates via $[v]_\gamma=[I_V]^{\gamma}_{\beta}[v]_\beta$.
\item \textbf{Change of representation:} use $[T]^{\gamma}_{\gamma}=[I_V]^{\gamma}_{\beta}[T]^{\beta}_{\beta}[I_V]^{\beta}_{\gamma}$.
\item \textbf{Similarity proofs:} show $B=Q^{-1}AQ$; for invariants use $\det(Q^{-1}AQ)=\det(A)$ and $\operatorname{tr}(Q^{-1}AQ)=\operatorname{tr}(A)$.
\item \textbf{Eigenvalues:} compute $p_T(\lambda)=\det([T]-\lambda I)$ and solve $p_T(\lambda)=0$ over the \emph{base field}.
\item \textbf{Eigenspaces:} for each eigenvalue $\lambda$, solve $(T-\lambda I)v=0$ (a kernel computation).
\end{itemize}
\end{examtip}

\begin{extension}[title={Extension (Tutorial 6, Question 5): existence of polynomial solutions via triangularity}]
\textbf{Syllabus link:} MH1201 Weeks 5--6 — Rank--Nullity and invertibility.

\textbf{Theorem.} Let $n\in\mathbb{N}_0$ and $Q\in\mathcal{P}_n(\mathbb{R})$. Then there exists $P\in\mathcal{P}_n(\mathbb{R})$ such that
\[
Q=(X^2+X)P''+2XP'+P(3).
\]

\textbf{Proof strategy.}
Define $L:\mathcal{P}_n(\mathbb{R})\to\mathcal{P}_n(\mathbb{R})$ by $L(P)=(X^2+X)P''+2XP'+P(3)$.
Compute $L(X^k)$ for each $k$: the leading term of $L(X^k)$ has degree $k$ with coefficient $k(k+1)$ (for $k\ge 1$) and $L(1)=1$.
Hence the matrix of $L$ w.r.t.\ the standard basis $(1,X,\dots,X^n)$ is \emph{lower triangular} with diagonal entries $1,2,6,12,\dots,n(n+1)$, all nonzero.
A triangular matrix with nonzero diagonal is invertible, so $L$ is an isomorphism and in particular surjective. Therefore, for every $Q$, the equation $L(P)=Q$ has a solution.

\textbf{Key technique:} When an operator on a polynomial space has a triangular matrix representation with nonzero diagonal, it is automatically invertible. This is a common pattern on MH1201 exams.
\end{extension}

\begin{extension}[title={Extension (Tutorial 6, Question 6): isomorphism between upper-triangular and symmetric matrices}]
\textbf{Syllabus link:} MH1201 Weeks 5--6 — Invertibility and isomorphisms.

\textbf{Theorem.} Let $V=\{M\in\mathbb{R}^{n\times n}\mid M\text{ upper triangular}\}$ and $W=\{M\in\mathbb{R}^{n\times n}\mid M\text{ symmetric}\}$. Then $V\cong W$.

\textbf{Proof.} Both $V$ and $W$ have dimension $\frac{n(n+1)}{2}$ (each is determined by the upper-triangular entries).
An explicit isomorphism $\Phi:V\to W$ is given by ``mirror the upper triangle'':
\[
\Phi(U)_{ij}\coloneqq \begin{cases}u_{ij},& i\le j,\\u_{ji},& i>j.\end{cases}
\]
This is linear (entrywise), and its inverse $\Psi:W\to V$ extracts the upper-triangular part: $\Psi(S)_{ij}=s_{ij}$ for $i\le j$ and $0$ otherwise. Both compositions are the identity on their respective spaces.
\end{extension}

\begin{examplebox}
\textbf{Worked Example (Tutorial 5, Question 4): substitution operator on polynomials.}
Let $Q\in\mathcal{P}(\mathbb{F})$ and define $T:\mathcal{P}_n(\mathbb{F})\to\mathcal{P}_{n\cdot\deg(Q)}(\mathbb{F})$ by $T(P)=P(Q)$.
Then $T$ is linear because
\[
T(P_1+P_2)=(P_1+P_2)(Q)=P_1(Q)+P_2(Q)=T(P_1)+T(P_2),\quad T(cP)=(cP)(Q)=cP(Q)=cT(P).
\]
To verify $T$ maps into $\mathcal{P}_{n\cdot\deg(Q)}(\mathbb{F})$: if $\deg(P)\le n$ and $P\neq 0$, then $\deg(P(Q))=\deg(P)\cdot\deg(Q)\le n\cdot\deg(Q)$ (using the identity $\deg(P(Q))=\deg(P)\cdot\deg(Q)$ for nonzero polynomials).
\end{examplebox}

\begin{examplebox}
\textbf{Worked Example (Tutorial 5, Question 2): linearity of the matrix-of-a-map operator.}
The map $M:\mathcal{L}(V,W)\to\mathbb{F}^{n\times m}$, $M(T)=[T]^{\beta}_{\alpha}$, is linear.
\textbf{Proof sketch:} Column $i$ of $M(T)$ is $[T(v_i)]_\beta$. Since the coordinate map $[\,\cdot\,]_\beta$ is linear,
\[
M(S+T)_{*,i}=[(S+T)(v_i)]_\beta=[S(v_i)]_\beta+[T(v_i)]_\beta=M(S)_{*,i}+M(T)_{*,i}.
\]
Similarly $M(aT)=aM(T)$. Hence $M$ is linear.
This result underlies the isomorphism $\mathcal{L}(V,W)\cong\mathbb{F}^{n\times m}$.
\end{examplebox}

%==================================================
% USER COMMENTS (EDITABLE)
%==================================================
% - (Write personal reminders, TODOs, or links here.)
% - (E.g., "Re-derive the Rank-Nullity proof before midterm", "Memorize the rank-nullity equivalences", etc.)
% - (Review Tutorial 5 Q3 for non-direct-sum examples.)
% - (Practice change-of-basis computations from Tutorial 6 Q7.)
% - (Remember: eigenvalues are field-specific!)
